\begin{helper}[(Q2) integrated along a segment: a growth bound on $A$ itself]
Let $n\ge2$ and $p>1$.  There is a constant $c=c(n,p)>0$ such that the
following holds.  Let $\varepsilon\ge0$, let $\gamma_0,\gamma_1>0$, let
$\Omega\subseteq\R^n$, and let $A\colon\R^n\times\R\times\R^n\to\R^n$ satisfy
\noderef{EllipticStructureQ12} with these data --- \emph{all} of that node's
clauses, as that node states them; they are not re-expanded here, and
$D(\varepsilon)=\{h\in\R^n:\varepsilon+|h|^2\ne0\}$ is the set on which its
differentiability clause and its two bounds are stated.
Then for a.e.\ $x\in\R^n$: if $x\in\Omega$, then for every $z\in\R$ and every
$h\in\R^n$,
\[
|A(x,z,h)|\ \le\ |A(x,z,0)|
\;+\;c\,\gamma_1\bigl(\varepsilon+|h|^2\bigr)^{\frac p2-1}|h| .
\]
All powers are real powers of nonnegative reals, with the convention
$0^{s}=0$ for $s\ne0$.

The constant must be allowed to depend on $p$ and not on $n$ alone: for
$1<p<2$ the weight is singular at the origin, and the bound carries the factor
$\frac1{p-1}$ produced by $\int_0^1t^{p-2}\,dt$.  Since
$|h|\le(\varepsilon+|h|^2)^{1/2}$, the right-hand side is at most
$|A(x,z,0)|+c\,\gamma_1(\varepsilon+|h|^2)^{\frac{p-1}2}$; for
$0\le\varepsilon\le1$ this is the thesis' remark on p.~2 that (Q4) of the thesis
follows from (Q2) of the thesis when $A(x,z,0)=0$, and it is the form from which
consumers deduce $A(\cdot,u,\nabla u)\in L^{p/(p-1)}$ out of $\nabla u\in L^p$.
\end{helper}
\begin{proof}
Write $s=\frac p2-1$ and set
\[
c_1=\max\Bigl\{1,\ \tfrac1{p-1}\Bigr\},\qquad c=n\,c_1 .
\]
Both are finite and positive: $p>1$ makes $\frac1{p-1}$ finite and positive,
and $n\ge2$ makes $n>0$.  The constant depends only on $n$ and $p$, which is
the required quantifier order.  Note that $\gamma_0$ and (Q1) of the thesis play no
part below; only (Q2) of the thesis and the $C^1$ clause of
\noderef{EllipticStructureQ12} are used.

\emph{Choice of $x$.}  By \noderef{EllipticStructureQ12} the condition
(Q2) of the thesis holds for a.e.\ $x$, i.e.\ there is a null set $N$ such that for
every $x\notin N$ with $x\in\Omega$, (Q2) of the thesis holds at $x$ for all
$z,h,i,j$.  Since the conclusion to be proved is itself an
``a.e.\ $x$'' statement, it suffices to prove it for each such $x$.  Fix
$x\notin N$ with $x\in\Omega$, and fix $z\in\R$ and $h\in\R^n$.  (The set
$\R^n\setminus N$ is nonempty, indeed of full measure, because Lebesgue measure
on $\R^n$ is not the zero measure.)

\emph{Step 0: the case $h=0$.}  If $h=0$ the claim reads
$|A(x,z,0)|\le|A(x,z,0)|+c\gamma_1(\varepsilon)^{s}\cdot0$, and the added term
is $0$; so the inequality holds.  Assume from now on $h\ne0$, so that $|h|>0$
and $\varepsilon+|h|^2>0$.

\emph{Step 1: the operator bound coming from (Q2) of the thesis.}  Let
$g\in D(\varepsilon)$ and write $D A(x,z,g)$ for the differential of
$A(x,z,\cdot)$ at $g$, which exists because $D(\varepsilon)$ is open and
$A(x,z,\cdot)$ is of class $C^1$ there by \noderef{EllipticStructureQ12}, since
$x\in\Omega$; and (Q2) of the thesis is available at $g$ for the same reason, $g$
lying in $D(\varepsilon)$.  Writing
$h=\sum_{j=1}^nh^je_j$ in the standard basis and using linearity of
$DA(x,z,g)$,
\[
\bigl(DA(x,z,g)h\bigr)^i=\sum_{j=1}^nh^j\,D_{h_j}A^i(x,z,g)
\qquad(i=1,\dots,n).
\]
By (Q2) of the thesis, $|D_{h_j}A^i(x,z,g)|\le\gamma_1(\varepsilon+|g|^2)^{s}$ for
all $i,j$, so
\[
\bigl|\bigl(DA(x,z,g)h\bigr)^i\bigr|
\ \le\ \gamma_1\bigl(\varepsilon+|g|^2\bigr)^{s}\sum_{j=1}^n|h^j|
\ \le\ \sqrt n\,\gamma_1\bigl(\varepsilon+|g|^2\bigr)^{s}\,|h| ,
\]
the last step being Cauchy--Schwarz (Preamble) applied to the vectors
$(|h^1|,\dots,|h^n|)$ and $(1,\dots,1)$.  Summing the squares over $i$ and
taking square roots,
\[
\bigl|DA(x,z,g)h\bigr|
=\Bigl(\sum_{i=1}^n\bigl|\bigl(DA(x,z,g)h\bigr)^i\bigr|^2\Bigr)^{1/2}
\ \le\ \sqrt n\cdot\sqrt n\,\gamma_1\bigl(\varepsilon+|g|^2\bigr)^{s}|h|
= n\,\gamma_1\bigl(\varepsilon+|g|^2\bigr)^{s}|h| .
\tag{$*$}
\]

\emph{Step 2: the fundamental theorem of calculus along $[0,h]$.}  Define
$\Psi\colon\R\to\R^n$ by $\Psi(t)=A(x,z,th)$.  First, $\Psi$ is continuous on
$[0,1]$: $t\mapsto th$ is continuous and $A(x,z,\cdot)$ is continuous on all of
$\R^n$ by the first clause of \noderef{EllipticStructureQ12}.  Second, for
$t\in(0,1]$ the point $th$ lies in $D(\varepsilon)$, because $h\ne0$ and $t>0$
give $\varepsilon+|th|^2\ge t^2|h|^2>0$; since $D(\varepsilon)$ is open and
$A(x,z,\cdot)$ is $C^1$ there, and $t\mapsto th$ is linear, the chain rule gives
that $\Psi$ is differentiable at every such $t$ with
\[
\Psi'(t)=DA(x,z,th)\,h\qquad(0<t\le1).
\]
At $t=0$ the point $0$ lies in $D(\varepsilon)$ only when $\varepsilon\ne0$, and
nothing below needs differentiability there: the single excluded point is an
\emph{endpoint} of the segment $[0,h]$, so the variant of the fundamental
theorem of calculus that asks for differentiability on the open interval and
continuity on the closed one (Preamble) is exactly the right tool.
We check that $\Psi'$ is interval-integrable on $[0,1]$.  It is measurable:
by the Preamble item on measurability of the derivative, the map
$g\mapsto DA(x,z,g)h$ is measurable, and $t\mapsto th$ is continuous, so the
composition $t\mapsto DA(x,z,th)h$ is measurable, hence strongly measurable
because $\R^n$ is separable.  (Only differentiability of $A$ in $h$ is used
here, which the $C^1$ clause of \noderef{EllipticStructureQ12} supplies a
fortiori; no continuity of $g\mapsto DA(x,z,g)$ is needed at this step, which
is why the Preamble item on measurability of an arbitrary derivative is the
right tool.)  It is
dominated: by $(*)$ with $g=th$, and because $|th|=|t|\,|h|$ gives
$|th|^2=t^2|h|^2$,
\[
|\Psi'(t)|\ \le\ n\,\gamma_1\,|h|\,\bigl(\varepsilon+t^2|h|^2\bigr)^{s}
\qquad(0<t\le1),
\]
hence for almost every $t\in[0,1]$, the excluded endpoint $t=0$ being a null
set; an almost-everywhere domination is all the Preamble item on domination
asks for.
The right-hand side is interval-integrable on $[0,1]$: it is the constant
$n\gamma_1|h|$ times $t\mapsto(\varepsilon+|th+(1-t)\cdot0|^2)^{s}$, which is
interval-integrable by \noderef{SegmentWeightIntegrable} applied with the
segment endpoints $a=h$ and $b=0$.  By the Preamble item on domination,
$\Psi'$ is interval-integrable on $[0,1]$.  Hence the open-interval form of the
fundamental theorem of calculus (Preamble) --- differentiability on $(0,1)$ from
the previous paragraph, continuity on $[0,1]$ from the first clause of
\noderef{EllipticStructureQ12}, interval-integrability of the derivative just
established --- applies and gives
\[
A(x,z,h)-A(x,z,0)=\Psi(1)-\Psi(0)=\int_0^1\Psi'(t)\,dt .
\]

\emph{Step 3: estimating the integral.}  By the norm inequality for the
integral and then monotonicity (both Preamble), using the domination of
Step 2,
\[
\bigl|A(x,z,h)-A(x,z,0)\bigr|\ \le\ \int_0^1|\Psi'(t)|\,dt
\ \le\ n\,\gamma_1\,|h|\int_0^1\bigl(\varepsilon+t^2|h|^2\bigr)^{s}dt .
\]
We claim
\[
\int_0^1\bigl(\varepsilon+t^2|h|^2\bigr)^{s}dt
\ \le\ c_1\bigl(\varepsilon+|h|^2\bigr)^{s} .
\tag{$**$}
\]
The integral on the left is finite by \noderef{SegmentWeightIntegrable}, as
noted in Step 2.

\emph{Proof of $(**)$ when $p\ge2$, i.e.\ $s\ge0$.}  For $t\in[0,1]$ we have
$t^2\le1$, hence $\varepsilon+t^2|h|^2\le\varepsilon+|h|^2$; since $s\ge0$ the
map $u\mapsto u^{s}$ is nondecreasing on $[0,\infty)$, so
$(\varepsilon+t^2|h|^2)^{s}\le(\varepsilon+|h|^2)^{s}$ for every $t\in[0,1]$.
Integrating over the unit interval gives $(**)$ with the constant $1\le c_1$.

\emph{Proof of $(**)$ when $1<p<2$, i.e.\ $-\tfrac12<s<0$.}  Let
$t\in(0,1]$.  Since $t^2\le1$ and $\varepsilon\ge0$ we have
$t^2\varepsilon\le\varepsilon$, hence
\[
0\ <\ t^2\bigl(\varepsilon+|h|^2\bigr)=t^2\varepsilon+t^2|h|^2
\ \le\ \varepsilon+t^2|h|^2 ,
\]
the strict positivity because $t>0$ and $\varepsilon+|h|^2>0$ (Step 0).  As
$s<0$, the map $u\mapsto u^{s}$ is decreasing on $(0,\infty)$, so
\[
\bigl(\varepsilon+t^2|h|^2\bigr)^{s}
\ \le\ \bigl(t^2(\varepsilon+|h|^2)\bigr)^{s}
= t^{2s}\bigl(\varepsilon+|h|^2\bigr)^{s} ,
\]
where the product rule for real powers is available because both factors
$t^2$ and $\varepsilon+|h|^2$ are strictly positive.  The single point $t=0$
is a null set, so the bound holds a.e.\ on $[0,1]$, and both sides are
interval-integrable there: the left by \noderef{SegmentWeightIntegrable}, the
right because $2s=p-2>-1$ (Preamble item on real powers).  Integrating,
\[
\int_0^1\bigl(\varepsilon+t^2|h|^2\bigr)^{s}dt
\ \le\ \bigl(\varepsilon+|h|^2\bigr)^{s}\int_0^1t^{2s}\,dt
=\frac{\bigl(\varepsilon+|h|^2\bigr)^{s}}{2s+1}
=\frac{\bigl(\varepsilon+|h|^2\bigr)^{s}}{p-1}
\ \le\ c_1\bigl(\varepsilon+|h|^2\bigr)^{s},
\]
using $\int_0^1t^{r}\,dt=\frac1{r+1}$ for $r>-1$ and $2s+1=p-1>0$, and
$(\varepsilon+|h|^2)^{s}\ge0$.  This proves $(**)$ in the second case, and the
factor $\frac1{p-1}$ is exactly the $p$-dependence announced in the statement.

\emph{Conclusion.}  Combining Step 3 with $(**)$,
\[
\bigl|A(x,z,h)-A(x,z,0)\bigr|
\ \le\ n\,c_1\,\gamma_1\bigl(\varepsilon+|h|^2\bigr)^{s}|h|
= c\,\gamma_1\bigl(\varepsilon+|h|^2\bigr)^{s}|h| ,
\]
and the triangle inequality $|A(x,z,h)|\le|A(x,z,0)|+|A(x,z,h)-A(x,z,0)|$
yields the assertion.  Since $x$ was an arbitrary point outside the null set
$N$, the assertion holds for a.e.\ $x$, as claimed.
\end{proof}
